\documentclass[10pt,a4paper]{amsart}
\usepackage[margin=19mm]{geometry}
\usepackage{amsmath,amssymb,mathtools}
\usepackage[scr=rsfs,cal=euler]{mathalfa}
\usepackage{xfrac}
\usepackage[unicode,hidelinks,bookmarks=false]{hyperref}
\newcommand{\ie}{i.e.\ }
\newcommand{\cf}{cf.\ }
\newcommand{\resp}{respectively\ }
\newcommand{\Subsection}{\subsection}
\providecommand{\nobreakdash}{\nobreak\hbox{-}}
\setlength{\emergencystretch}{2em}
\multlinegap0pt
\usepackage{bm}
\usepackage{bbm}
\usepackage{dsfont}
\usepackage{xspace}
\usepackage{cancel}
\usepackage{mathtools}
\usepackage{amsthm}
\makeatletter
\@ifundefined{theorem}{\newtheorem{theorem}{Theorem}}{}
\@ifundefined{lemma}{\newtheorem{lemma}[theorem]{Lemma}}{}
\makeatother
\newcommand{\abs}[1]{\lvert #1 \rvert}
\title[A generalized Helmholtz-type decomposition of symmetric tens]{A generalized Helmholtz-type decomposition of symmetric tensor fields and applications to ray transforms}
\author{Antti Kykk\"anen, Rohit Kumar Mishra, Suman Kumar Sahoo}
\date{Source: arXiv:2602.18983v2 (CC BY 4.0)}
\begin{document}
\maketitle

\noindent\emph{Source and licence.} A generalized Helmholtz-type decomposition of symmetric tensor fields and applications to ray transforms. Version arXiv:2602.18983v2 (CC BY 4.0), distributed under the Creative Commons Attribution 4.0 International licence, as recorded in the arXiv licence metadata for this exact version and shown on the abstract page https://arxiv.org/abs/2602.18983v2. Full rights notice: licenses/arxiv-2602.18983-CC-BY-4.0.txt.\par\smallskip

\noindent\textit{Editorial scope.} Counted unit: the Lemma whose source label is \texttt{lem:E2-decomposition} in the article (source0.tex lines 737--778, source label \texttt{lem:E2-decomposition}), delivered together with the article's own complete original proof of it (source0.tex lines 780--1310, one \texttt{proof} environment of 531 lines). No mathematics is written, repaired or re-derived: statement and proof are the article's own text line for line. Required context retained uncounted: none. Every definition, display and label of those lines is retained unchanged. Omitted from this package and not redistributed: 1--736, 779--779, 1311--2265. Nothing inside a retained excerpt is omitted or altered: every retained body is delivered exactly as the article's lines are written, so no omission receipt of any kind accompanies this package. Citations: the retained excerpts carry no citation command at all, so no bibliography entry of the article is transcribed and no reference is redistributed. Reference adaptation: none was needed and none was made; every reference inside the delivered unit resolves inside it, so no unresolved reference or citation is produced. Printed slips kept literal: no printed word, symbol, hypothesis or number of the counted statement or of its proof was altered. Layout and class: the article's own class is \texttt{article} with packages including hyperref, amsthm, amssymb, amsfonts, amsmath, amsrefs, amscd, color, fancyhdr, graphicx, wrapfig, lipsum, soul, enumitem; the wrapper is the lane's pinned \texttt{amsart} editorial wrapper at 10pt on A4, which restates the theorem-like environments the retained text needs and adds the article's own macro definitions that the retained text needs (\texttt{\textbackslash abs}), with extra wrapper packages (bm, bbm, dsfont, xspace, cancel, mathtools). The delivered numbering is the unit's own. Assets: no retained span contains a figure environment, a graphics-inclusion command, a PDF-page inclusion, a table or a TikZ picture; no image file is copied into this package, the package declares no asset dependency and no image file is redistributed with it. Licence: version arXiv:2602.18983v2 of this article is distributed under the Creative Commons Attribution 4.0 International licence, as recorded in the arXiv licence metadata for this exact version and shown on the abstract page https://arxiv.org/abs/2602.18983v2. A full rights notice is delivered at licenses/arxiv-2602.18983-CC-BY-4.0.txt. Characters: every retained excerpt is pure ASCII, so the delivered engine, XeLaTeX with its default Latin Modern text font, reports no missing character.
\begin{lemma}[Decomposition of elastic $2$-tensors]
\label{lem:E2-decomposition}
Let $y \in \mathbb{R}^n\setminus \{0\}$. For any elastic tensor $\widehat f \in E^2(n)$ there are unique $\widehat u \in \mathbb{R}^n$, $\widehat v \in E^1(n)$ and $\widehat g \in E^2(n)$ so that $\widehat f = \widehat{H}_y \widehat v + \widehat{K}_y \widehat u + \widehat g$ with $\widehat{H}^*_y\widehat g = \widehat{K}^*_y\widehat g = 0$ and $\widehat u \perp y$.

Moreover, $\widehat v$ and $\widehat u$ can be explicitly computed in terms of $\widehat f$ by
\begin{equation}
\label{eqn:hat-u}
\widehat u_i
=
\frac{4}{n-1}
\left(
\frac{1}{\abs{y}^2}
\sum_j
\varepsilon_{ij}(y)
(\widehat{K}^*_y\widehat f)_j
-
\frac{1}{\abs{y}^4}
\sum_k
y_k(\widehat{H}^*_y\widehat f)_{ik}
+
\frac{y_i}{\abs{y}^6}
\sum_{i,k}
y_iy_k
(\widehat{H}^*_y\widehat f)_{ik}
\right)
\end{equation}
and
\begin{equation}
\label{eqn:hat-v}
\widehat v_{ij}
=
\frac{2}{\abs{y}^4}
(\widehat{H}^*_y\widehat f)_{ij}
-
\frac{y_iy_j}{\abs{y}^8}
\sum_{k,l}y_ky_l
(\widehat{H}^*_y\widehat f)_{kl}
-
\frac{1}{2\abs{y}^2}(y_j\widehat u_i + y_i\widehat u_j)
\end{equation}
where $\varepsilon_{ij}(y) = \delta_{ij}-\frac{y_iy_j}{\abs{y}^{2}}$.
\end{lemma}

\begin{proof}
We start by proving uniqueness of such a decomposition. Assume that $\widehat{f} \in E^2(n)$, $ \widehat g \in E^2(n)$, $\widehat v \in E^1(n)$ and $\widehat u \in \mathbb{R}^n$ are such that
\begin{equation}
\label{eqn:assumed-decomposition}
\widehat f
=
\widehat{H}_y\widehat v
+
\widehat{K}_y\widehat u
+
\widehat g
\quad\text{with}\quad
\widehat{H}^*_y\widehat g
=
\widehat{K}^*_y\widehat g
=
0
\quad\text{and}\quad
\widehat u \perp y.
\end{equation}
Uniqueness of the decomposition is proved by showing that given this decomposition of $\widehat f$ then $\widehat u$ and $\widehat v$ are necessarily of the form~\eqref{eqn:hat-u} and~\eqref{eqn:hat-v}.

To show that~\eqref{eqn:hat-u} and~\eqref{eqn:hat-v} we start by applying the operator $\widehat{K}^*_y$ to~\eqref{eqn:assumed-decomposition} to obtain
\begin{equation}
\label{eqn:k-star-applied-to-decomp}
\widehat{K}^*_y\widehat f
=
\widehat{K}^*_y\widehat{H}_y\widehat v
+
\widehat{K}^*_y\widehat{K}_y\widehat u
\end{equation}
since $\widehat{H}^*_y\widehat g = \widehat{K}^*_y\widehat g = 0$. We compute that in components
\begin{equation}
\label{eqn:k-star-f-decomp-part-1}
\begin{split}
(\widehat{K}^*_y\widehat{H}_y\widehat v)_i
&=
\sum_{k,j}
y_j(\widehat{H}_y\widehat v)_{ijkk}=
\frac{1}{2}
\sum_{k,j}
\left(
y_iy_j^2\widehat v_{kk}
+
y_jy_k^2\widehat v_{ij}
\right)
\\
&=
\frac{\abs{y}^2}{2}
\left(
\sum_{k}
y_i\widehat v_{kk}
+
\sum_{j}
y_j\widehat v_{ij}
\right)=
\frac{\abs{y}^2}{2}
\sum_{k}
\left(
y_i\widehat v_{kk}
+
y_k\widehat v_{ik}
\right).
\end{split}
\end{equation}
For the second part of the right-hand side of~\eqref{eqn:k-star-applied-to-decomp}, we get
\begin{equation}
\begin{split}
(\widehat{K}^*_y\widehat{H}_y\widehat u)_i
&=
\sum_{k,j}
y_j(\widehat{K}_y\widehat u)_{ijkk}=
\frac{1}{4}
\sum_{k,j}
\left(
y_j(y_i\widehat u_j + y_j\widehat u_i)\delta_{kk}
+
y_j(y_k\widehat u_k+y_k\widehat u_k)\delta_{ij}
\right).
\end{split}
\end{equation}
Using the fact that $\widehat u \perp y$, the above reduces to
\begin{equation}
\label{eqn:k-star-f-decomp-part-2}
\begin{split}
(\widehat{K}^*_y\widehat{H}_y\widehat u)_i
&=
\frac14
\sum_{k,j}
y^2_j\widehat u_i\delta_{kk}=
\frac{n}{4}
\abs{y}^2
\widehat u_i.
\end{split}
\end{equation}
Combining~\eqref{eqn:k-star-f-decomp-part-1} and~\eqref{eqn:k-star-f-decomp-part-2} gives
\begin{equation}
\label{eqn:hat-k-f}
(\widehat{K}^*_y\widehat f)_i
=
\frac{n}{4}
\abs{y}^2
\widehat u_i
+
\frac{\abs{y}^2}{2}
\sum_{k}
\left(
y_i\widehat v_{kk}
+
y_k\widehat v_{ik}
\right).
\end{equation}
We will make use of this formula later, but first we will derive a formula for $\widehat{H}^*_y\widehat f$. We compute in coordinates that
\begin{equation}
\label{eqn:h-star-f-decomp-part-1}
\begin{split}
(\widehat{H}^*_y\widehat{H}_y\widehat v)_{ij}
&=
\sum_{k,l}
y_ky_l(\widehat{H}_y\widehat v)_{ijkl}
=
\frac12
\sum_{k,l}
y_ky_l
\left(
y_iy_j\widehat v_{kl}
+
y_ky_l\widehat v_{ij}
\right)
=
\frac12
\left(
y_iy_j
\sum_{k,l}
y_ky_l\widehat v_{kl}
+
\abs{y}^4
\widehat v_{ij}
\right)
.
\end{split}
\end{equation}
Similarly, we compute that
\begin{equation}
\begin{split}
(\widehat{H}^*_y\widehat{K}_y\widehat u)_{ij}
&=
\sum_{k,l}
y_ky_l
(\widehat{K}_y\widehat u)_{ijkl}
=
\frac{1}{4}
\sum_{k,l}
y_ky_l
\left(
(y_i\widehat u_j + y_j\widehat u_i)\delta_{kl}
+
(y_k\widehat u_l + y_l\widehat u_k)\delta_{ij}
\right)
\\
&=
\frac{1}{4}
\left(
\abs{y}^2
(y_j\widehat u_i + y_i\widehat u_i)
+
2
\abs{y}^2
(\widehat u \cdot y)
\delta_{ij}
\right).
\end{split}
\end{equation}
Applying the fact that $\widehat u \perp y$, reduces the above to
\begin{equation}
\label{eqn:h-star-f-decomp-part-2}
\begin{split}
(\widehat{H}^*_y\widehat{K}_y\widehat u)_{ij}
=
\frac{\abs{y}^2}{4}
\left(
y_j\widehat u_i+y_i\widehat u_j
\right).
\end{split}
\end{equation}
Since $\widehat{H}^*_y\widehat g = \widehat{K}^*_y\widehat g = 0$,  combining~\eqref{eqn:h-star-f-decomp-part-1} and~\eqref{eqn:h-star-f-decomp-part-2} gives
\begin{equation}
\label{eqn:h-star-f}
(\widehat{H}^*_y\widehat f)_{ij}
=
\frac{1}{2}
\left(
y_iy_j\sum_{k,l}
y_ky_l\widehat v_{kl}
+
\abs{y}^4
\widehat v_{ij}
\right)
+
\frac{\abs{y}^2}{4}
\left(
y_j\widehat u_i+y_i\widehat u_j
\right).
\end{equation}
We contract this equation with $y$ twice (once in both the $i$ and $j$ indices) and use the definition of $\widehat{H}^*_y$ to obtain
\begin{equation}
\label{eqn:f-contracted-al-indices}
\begin{split}
\sum_{i,j,k,l}
y_iy_jy_ky_l
\widehat f_{ijkl}
&=
\sum_{i,j}
y_iy_j(\widehat{H}^*_y\widehat f)_{ij}
\\
&=
\frac12
\left(
\sum_i
y_i^2
\right)
\left(
\sum_i
y_i^2
\right)
\left(
\sum_{k,l}
y_ky_l\widehat v_{kl}
\right)
+
\frac12\abs{y}^4
\sum_{i,j}
y_iy_j\widehat v_{ij}
\\
&\quad
+
\frac{\abs{y}^2}{2}
\left(
\left(
\sum_j y_j^2
\right)
\left(
\sum_i
y_i\widehat u_i
\right)
+
\left(
\sum_i y_i^2
\right)
\left(
\sum_k
y_j\widehat u_j
\right)
\right)
\\
&=
\abs{y}^4
\sum_{k,l}y_ky_l\widehat v_{kl}
+
\abs{y}^4(\widehat u \cdot y)=
\abs{y}^4
\sum_{k,l}y_ky_l\widehat v_{kl}.
\end{split}
\end{equation}
We now go back to~\eqref{eqn:k-star-f-decomp-part-1} and~\eqref{eqn:k-star-f-decomp-part-2} and use them to compute
\begin{equation}
\begin{split}
\sum_{i}
y_i(\widehat{K}_y^*\widehat f)_i
&=
\sum_{i}
y_i(\widehat{K}_y^*\widehat{H}_y\widehat v)_i
+
\sum_{i}
y_i(\widehat{K}_y^*\widehat{K}_y\widehat u)_i
\\
&=
\frac{\abs{y}^2}{2}
\sum_{i}
y_i
\left(
y_i
\sum_{k}
\widehat v_{kk}
+
\sum_{k}
y_k\widehat v_{ik}
\right)
+
\sum_{i}
y_i
\left(
\frac{n}{2}
\abs{y}^2\widehat u_i
\right)
\\
&=
\frac{\abs{y}^4}{2}
\sum_{k}
\widehat v_{kk}
+
\frac{\abs{y}^2}{2}
\sum_{i,k}
y_iy_k
\widehat v_{ik}
+
\frac{n}{2}\abs{y}^2(\widehat u \cdot y)
\\
&=
\frac{\abs{y}^4}{2}
\sum_{k}
\widehat v_{kk}
+
\frac{\abs{y}^2}{2}
\sum_{i,k}
y_iy_k
\widehat v_{ik}.
\end{split}
\end{equation}
From the above equation, we compute using~\eqref{eqn:f-contracted-al-indices} that
\begin{equation}
\begin{split}
\sum_k
\widehat v_{kk}
&=
\frac{2}{\abs{y}^4}
\sum_i
y_i(\widehat{K}^*_y\widehat f)_i
-
\frac{1}{2\abs{y}^2}
\sum_{i,k}
y_iy_k\widehat v_{ik}
=
\frac{1}{\abs{y}^4}
\left(
2
\sum_i
y_i(\widehat{K}^*_y\widehat f)_i
-
\frac{1}{\abs{y}^2}
\sum_{i,j,k,l}
y_iy_jy_ky_l
\widehat f_{ijkl}
\right).
\end{split}
\end{equation}
We insert this into~\eqref{eqn:hat-k-f} to get
\begin{equation}
\label{eqn:auxiliary-equation-1}
\begin{split}
\frac{n}{4}\abs{y}^2\widehat u_i
+
\frac{\abs{y}^2}{2}
\sum_k y_k\widehat v_{ik}
&=
(\widehat{K}^*_y\widehat f)_i
-
\frac{\abs{y}^2}{2}
y_i
\sum_k \widehat v_{kk}
\\
&=
(\widehat{K}^*_y\widehat f)_i
-
\frac{y_i}{2\abs{y}^2}
2\sum_{j}y_j(\widehat{K}^*_y\widehat f)_j
-
\frac{y_i}{2\abs{y}^4}
\sum_{i,j,k,l}
y_iy_jy_ky_l
\widehat f_{ijkl}
.
\end{split}
\end{equation}
Then we go back to~\eqref{eqn:h-star-f} to compute
\begin{equation}
\begin{split}
\sum_{j,k,l}y_jy_ky_l\widehat f_{ijkl}
&=
\sum_{j}
y_j(\widehat{H}^*_y\widehat f)_{ij}
\\
&=
\sum_jy_j
\left(
\frac12
\left(
y_iy_j
\sum_{k,l}
y_ky_l\widehat v_{kl}
+
\abs{y}^4
\widehat v_{ij}
\right)
+
\frac{\abs{y}^4}{4}
\left(
y_j\widehat u_i
+
y_i\widehat u_j
\right)
\right)
\\
&=
\frac{\abs{y}^2}{2}
y_i
\sum_{k,l}
y_ky_l
\widehat v_{kl}
+
\frac{\abs{y}^4}{2}
\sum_jy_j\widehat v_{ij}
+
\frac{\abs{y}^4}{4}\widehat u_i
+
\frac{\abs{y}^4}{4}y_i(\widehat u \cdot y)
\\
&=
\frac{\abs{y}^2}{2}
y_i
\sum_{k,l}
y_ky_l
\widehat v_{kl}
+
\frac{\abs{y}^4}{2}
\sum_jy_j\widehat v_{ij}
+
\frac{\abs{y}^4}{4}\widehat u_i.
\end{split}
\end{equation}
Rearranging the above equation and inserting~\eqref{eqn:f-contracted-al-indices} gives
\begin{equation}
\label{eqn:auxiliary-equation-2}
\frac{\abs{y}^4}{4}\widehat u_i
+
\frac{\abs{y}^4}{2}
\sum_jy_j\widehat v_{ij}
=
\sum_{j,k,l}
y_jy_ky_l
\widehat f_{ijkl}
-
\frac{y_i}{2\abs{y}^2}
\sum_{i,j,k,l}
y_iy_jy_ky_l
\widehat f_{ijkl}.
\end{equation}
Then multiplying~\eqref{eqn:auxiliary-equation-1} by $\abs{y}^2$ and then subtracting~\eqref{eqn:auxiliary-equation-2} gives
\begin{equation}
\begin{split}
\frac{n-1}{4}\abs{y}^4\widehat u_i
&=
\abs{y}^2
\left(
(\widehat{K}^*_y\widehat f)_i
-
\frac{y_i}{\abs{y}^2}\sum_m y_m(\widehat{K}^*_y\widehat f)_m
+
\frac{y_i}{2\abs{y}^4}
\sum_{i,j,k,l}
y_iy_jy_ky_l\widehat f_{ijkl}
\right)
\\
&\quad
-
\sum_{j,k,l}
y_jy_ky_l\widehat f_{ijkl}
+
\frac{y_i}{2\abs{y}^2}
\sum_{i,j,k,l}
y_iy_jy_ky_l\widehat f_{ijkl}
\\
&=
\abs{y}^2
\sum_j
\varepsilon_{ij}(y)
(\widehat{K}^*_y\widehat f)_j
-
\sum_k
y_k(\widehat{H}^*_y\widehat f)_{ik}
+
\frac{y_i}{\abs{y}^2}
\sum_{i,k}
y_iy_k
(\widehat{H}^*_y\widehat f)_{ik},
\end{split}
\end{equation}
where $ \varepsilon_{ij}(y)
=
\delta_{ij}
-\frac{y_iy_j}{\abs{y}^2}.$
% \begin{equation}
% \varepsilon_{ij}(y)
% \coloneqq
% \delta_{ij}
% -
% \frac{y_iy_j}{\abs{y}^2}.
% \end{equation}
This proves formula~\eqref{eqn:hat-u} since the exact form can be obtained by diving by a power of $\abs{y}$ appropriately. Next, we will find a formula for $\widehat v$ in terms of $\widehat f$ and $\widehat u$ (and thus in terms of $\widehat f$ alone by~\eqref{eqn:hat-u}). In order to find the formula for $\widehat v$ rewrite~\eqref{eqn:h-star-f} and use~\eqref{eqn:f-contracted-al-indices} to compute
\begin{equation}
\begin{split}
\widehat v_{ij}
&=
\frac{1}{\abs{y}^4}
\left(
2(\widehat{H}^*\widehat f)_{ij}
-
y_iy_j
\sum_{k,l}
y_ky_l
\widehat v_{kl}
-
\frac{\abs{y}^2}{2}(y_j\widehat u_i+y_i\widehat u_j)
\right)
\\&=
\frac{1}{\abs{y}^4}
\left(
2(\widehat{H}^*\widehat f)_{ij}
-
\frac{y_iy_j}{\abs{y}^4}
\sum_{k,l}
(\widehat{H}^*_y\widehat f)_{kl}
-
\frac{\abs{y}^2}{2}(y_j\widehat u_i+y_i\widehat u_j)
\right).
\end{split}
\end{equation}
This proves the second formula (Equation~\eqref{eqn:hat-v}) of the lemma. The exact form is obtained by distributing the powers of $\abs{y}$ above.

Finally, to find $\widehat g$ in terms of $\widehat f$ we can simply substract $\widehat{H}_y\widehat v + \widehat{K}_y\widehat u$ from $\widehat f$ since $\widehat u$ and $\widehat v$ can be written in terms of $\widehat f$. Given the desired decomposition, we have derived formulas for $\widehat v$, $\widehat u$ and $\widehat g$ in terms of $\widehat f$. This shows that such a decomposition is necessarily unique. To prove existence, we can simply define $\widehat v$, $\widehat u$ and $\widehat g$ using the formulas above. This completes the proof of our lemma.
\end{proof}

\end{document}
